\documentclass[11pt]{article}
\usepackage[left=1in,right=1in,top=0.9in,bottom=0.9in]{geometry}
\usepackage{amsmath,amssymb}
\usepackage{booktabs}
\usepackage[round]{natbib}
\usepackage{hyperref}
\hypersetup{colorlinks,citecolor=blue,filecolor=blue,linkcolor=blue,urlcolor=blue}
\usepackage{float}
\usepackage[ruled,lined]{algorithm2e}
\usepackage{microtype}
\usepackage{placeins}
\usepackage{titling}
\setlength{\droptitle}{-2em}
\posttitle{\par\end{center}\vspace{-1em}}

\newcommand{\E}{\mathbb{E}}
\newcommand{\R}{\mathbb{R}}
\newcommand{\calA}{\mathcal{A}}
\newcommand{\calS}{\mathcal{S}}
\newcommand{\calH}{\mathcal{H}}

\title{Primer: Maximum Causal Entropy IRL (MCE-IRL)}
\author{econirl package documentation}
\date{}

\begin{document}
\maketitle


%% ============================================================
\section{Problem and Motivation}
%% ============================================================

Standard maximum entropy IRL \citep{ziebart2008} recovers a reward function from demonstrations by matching feature expectations under the maximum entropy principle. However, it uses non-causal state visitation frequencies that propagate future information backward, violating the agent's information structure. When features vary by action, this produces incorrect policies.

\citet{ziebart2010} introduced Maximum Causal Entropy IRL (MCE-IRL), which conditions each action only on information available at decision time. The causal entropy $H(A_t | S_{1:t}, A_{1:t-1})$ replaces the Shannon entropy $H(A^T)$, ensuring that the recovered policy respects the sequential timing of decisions. This is exactly what economists call ``rational expectations'' in dynamic programming.

The resulting softmax Bellman equation is simultaneously the MCE optimal policy (Theorem 1 of \citealt{ziebart2010}) and the structural logit choice model of \citet{rust1987}. An economist sees structural estimation with feature matching. A machine learning researcher sees maximum entropy IRL. They solve the same optimization problem. This makes MCE-IRL the bridge estimator between the two traditions.


%% ============================================================
\section{Model and Notation}
%% ============================================================

The agent solves the maximum causal entropy problem
\begin{equation}
\label{eq:mce_objective}
\max_\pi \; \E_\pi\!\left[\sum_{t=0}^{\infty} \gamma^t \bigl(r(s_t, a_t;\theta) + \sigma \calH(\pi(\cdot | s_t))\bigr)\right]
\end{equation}
subject to the feature matching constraint $\E_\pi[\phi] = \E_{\text{data}}[\phi]$, where $r(s,a;\theta) = \theta^\top \phi(s,a)$ is linear in known features $\phi$ with unknown weights $\theta$.

\begin{table}[H]
\centering\small
\begin{tabular*}{\textwidth}{@{\extracolsep{\fill}} l l l}
\toprule
Symbol & Name & Definition \\
\midrule
$Q(s,a)$ & Soft Q-value & $r(s,a;\theta) + \beta \sum_{s'} P(s'|s,a)\, V(s')$ \\
$V(s)$ & Soft value & $\sigma \log \sum_a \exp(Q(s,a)/\sigma)$ \\
$\pi(a|s)$ & Policy (CCP) & $\exp\bigl((Q(s,a) - V(s))/\sigma\bigr)$ \\
$D_\pi(s)$ & State visitation & Expected discounted visits to $s$ under $\pi$ \\
$\rho_0$ & Initial distribution & Empirical start-state frequencies \\
$F_\pi$ & Policy transitions & $F_\pi[s,s'] = \sum_a \pi(a|s)\, P(s'|s,a)$ \\
$\phi(s,a)$ & Features & Known, shape $(|\calS|, |\calA|, K)$ \\
$\theta$ & Reward weights & Unknown, dimension $K$ \\
\bottomrule
\end{tabular*}
\end{table}


%% ============================================================
\section{Derivations}
%% ============================================================

\subsection{Theorem 1: The MCE Optimal Policy}

We seek the policy that is maximally noncommittal (highest entropy) while reproducing the observed feature statistics. The \emph{causal} entropy conditions each action only on information available at decision time, respecting the agent's sequential information structure.

\paragraph{Primal problem.}
\begin{equation}
\label{eq:primal}
\max_{\pi} \; H_{\text{causal}}(\pi) \quad \text{s.t.} \quad \E_\pi[\phi] = \bar\phi_{\text{data}}, \quad \pi(a|s) \geq 0, \quad \sum_a \pi(a|s) = 1 \;\forall s.
\end{equation}
The causal entropy is
\begin{equation}
H_{\text{causal}}(\pi) = \sum_{t=0}^{\infty} \gamma^t \sum_{s} d_\pi^t(s) \sum_{a} \pi(a|s) \log \frac{1}{\pi(a|s)},
\end{equation}
where $d_\pi^t(s) = P(S_t = s \mid \pi)$ is the probability of being in state $s$ at time $t$ under policy $\pi$.

\paragraph{Lagrangian.}
Introducing multipliers $\theta \in \R^K$ for the feature constraints and $\lambda(s)$ for the normalization constraints,
\begin{equation}
\calL(\pi, \theta, \lambda) = H_{\text{causal}}(\pi) + \theta^\top\!\left(\E_\pi[\phi] - \bar\phi_{\text{data}}\right) + \sum_s \lambda(s)\!\left(\sum_a \pi(a|s) - 1\right).
\end{equation}
Substituting the definitions and taking the derivative with respect to $\pi(a|s)$,
\begin{equation}
\frac{\partial \calL}{\partial \pi(a|s)} = d_\pi(s)\!\left[-\log \pi(a|s) - 1 + \theta^\top \phi(s,a) + \text{continuation}\right] + \lambda(s) = 0.
\end{equation}
The ``continuation'' term accounts for how changing $\pi(a|s)$ affects future state visitation. Solving and absorbing constants into $\lambda(s)$, we obtain
\begin{equation}
\pi^*(a|s) = \frac{\exp\!\bigl(Q^*(s,a)/\sigma\bigr)}{\sum_{a'} \exp\!\bigl(Q^*(s,a')/\sigma\bigr)},
\end{equation}
where $Q^*$ satisfies the soft Bellman recursion
\begin{align}
Q(s,a) &= \theta^\top \phi(s,a) + \beta \sum_{s'} P(s'|s,a)\, V(s'), \label{eq:soft_q} \\
V(s) &= \sigma \log \sum_{a} \exp\bigl(Q(s,a)/\sigma\bigr). \label{eq:soft_v}
\end{align}
The policy is therefore
\begin{equation}
\label{eq:soft_pi}
\pi(a|s) = \exp\!\bigl((Q(s,a) - V(s))/\sigma\bigr).
\end{equation}
This is exactly the structural logit model of \citet{rust1987} with Type I extreme value errors at scale $\sigma$. The Lagrange multipliers $\theta$ play the role of the structural parameters, and the soft Bellman recursion is the standard DDC value function.


\subsection{Theorem 2: Feature Matching Gradient}

\paragraph{Dual problem.}
Define the dual function $G(\theta) = \max_\pi \calL(\pi, \theta)$. Since the inner maximization has the closed-form solution $\pi^*(\theta)$ from Theorem 1, we can write
\begin{equation}
G(\theta) = H_{\text{causal}}(\pi^*(\theta)) + \theta^\top\!\left(\E_{\pi^*}[\phi] - \bar\phi_{\text{data}}\right).
\end{equation}
The MCE-IRL problem is $\min_\theta G(\theta)$.

\paragraph{Gradient by the envelope theorem.}
Since $\pi^*$ is the unconstrained maximizer of $\calL$ for fixed $\theta$, the envelope theorem gives
\begin{equation}
\label{eq:gradient}
\nabla_\theta G(\theta) = \E_{\pi^*(\theta)}[\phi] - \bar\phi_{\text{data}} = \bar\phi_\pi(\theta) - \bar\phi_{\text{data}}.
\end{equation}
The expected features under the current policy are computed via the occupancy measure
\begin{equation}
\bar\phi_\pi(\theta) = \sum_{s} D_\pi(s) \sum_{a} \pi(a|s)\, \phi(s,a),
\end{equation}
where $D_\pi(s)$ is the expected discounted state visitation frequency under $\pi$. This satisfies the fixed-point equation
\begin{equation}
\label{eq:occupancy}
D_\pi = \rho_0 + \beta\, F_\pi^\top\, D_\pi \quad \Longrightarrow \quad D_\pi = (I - \beta\, F_\pi^\top)^{-1}\, \rho_0,
\end{equation}
with $F_\pi[s,s'] = \sum_a \pi(a|s)\, P(s'|s,a)$ and $\rho_0$ the initial state distribution.

At the optimum, $\nabla_\theta G = 0$, so $\bar\phi_\pi = \bar\phi_{\text{data}}$. The gradient vanishes when the policy's feature expectations match the data.

\paragraph{Equivalence to NFXP.}
The conditional log-likelihood of the DDC logit model is $\ell(\theta) = \sum_i \log \pi_\theta(a_i|s_i)$. Its gradient is
\begin{equation}
\nabla_\theta \ell(\theta) = \frac{1}{\sigma}\!\left(\bar\phi_{\text{data}} - \bar\phi_\pi(\theta)\right),
\end{equation}
which differs from Equation~\eqref{eq:gradient} only by the factor $1/\sigma$. MCE-IRL and NFXP optimize the same objective and converge to the same $\hat\theta$.


\subsection{Theorem 3: Robustness Guarantee}

The MCE distribution $P_{\text{MCE}}$ solves a minimax problem over the set $\calC$ of distributions matching the feature constraints.

\paragraph{Statement.}
Among all distributions in $\calC$, the MCE distribution minimizes the worst-case KL divergence to the true (unknown) distribution $P'$:
\begin{equation}
P_{\text{MCE}} = \argmin_{P \in \calC} \; \max_{P' \in \calC} \; D_{\text{KL}}(P' \| P).
\end{equation}

\paragraph{Proof sketch.}
The KL divergence $D_{\text{KL}}(P' \| P) = \E_{P'}[\log P'/P]$ is minimized (in $P$) when $P$ places as much probability mass as possible on all outcomes consistent with the constraints. The maximum entropy distribution does exactly this: it is the ``flattest'' distribution in $\calC$, making it the centroid of the constraint set in the KL geometry. Any other distribution $P \in \calC$ concentrates mass away from some region, creating a worst-case $P'$ that concentrates mass precisely in that region, yielding higher KL divergence. For a formal proof, see Chapter 3 of \citet{ziebart2010}.

\paragraph{Practical consequence.}
This means the MCE policy hedges against adversarial changes in the state distribution. NFXP, which solves the same optimization, inherits this guarantee. However, the guarantee applies only within the constraint set defined by the chosen features $\phi$. Misspecified features can still produce poor out-of-sample predictions.


\subsection{Identification}

Rewards are identified only up to potential-based shaping \citep{ng1999}. For any $h: \calS \to \R$, the shaped reward $r'(s,a) = r(s,a) + h(s) - \beta \sum_{s'} P(s'|s,a) h(s')$ produces the same optimal policy \citep{magnacThesmar2002}. The parametric restriction $r = \theta^\top \phi$ with $K < |\calS|(|\calA|-1)$ resolves this by constraining the reward to a low-dimensional subspace. MCE-IRL additionally constrains $\|\theta\|$ via the feature matching dual, selecting the maximum-entropy-compatible reward from the equivalence class \citep{kim2021}.


%% ============================================================
\section{Estimation Algorithm}
%% ============================================================

\begin{algorithm}[H]
\DontPrintSemicolon
\caption{MCE-IRL \citep{ziebart2010}}
\label{alg:mce_irl}
\KwIn{Panel $\{(s_i, a_i)\}_{i=1}^N$, features $\phi(s,a)$, transitions $P(s'|s,a)$, discount $\beta$, scale $\sigma$}
\BlankLine

\tcp{Step 1: Compute empirical feature expectations}
$\bar\phi_{\text{data}} = \frac{1}{N} \sum_{i=1}^N \phi(s_i, a_i)$\;
Compute empirical state occupancy $D_{\text{data}}(s) = \frac{1}{N}\sum_i \mathbf{1}[s_i = s]$\;
Compute initial state distribution $\rho_0$ from first observations\;
\BlankLine

\tcp{Step 2: Initialize reward weights}
$\theta_0 \leftarrow \mathbf{0}$\;
\BlankLine

\tcp{Step 3: Iterate until convergence}
\For{$k = 0, 1, \ldots$}{
    \BlankLine
    \tcp{3a. Backward pass: soft value iteration (Eqs.~\ref{eq:soft_q}--\ref{eq:soft_pi})}
    $R \leftarrow \theta_k^\top \phi$ \tcp*{reward matrix $(|\calS|, |\calA|)$}
    Solve $V = T_\sigma(V; R)$ via hybrid iteration to get $V$, $Q$, $\pi$\;
    \BlankLine
    \tcp{3b. Forward pass: occupancy measure (Eq.~\ref{eq:occupancy})}
    $F_\pi[s,s'] = \sum_a \pi(a|s)\, P(s'|s,a)$\;
    $D_\pi = (I - \beta\, F_\pi^\top)^{-1}\, \rho_0$\;
    \BlankLine
    \tcp{3c. Expected features under current policy}
    $\bar\phi_\pi = \sum_s D_\pi(s) \sum_a \pi(a|s)\, \phi(s,a)$\;
    \BlankLine
    \tcp{3d. Feature matching gradient (Eq.~\ref{eq:gradient})}
    $g = \bar\phi_{\text{data}} - \bar\phi_\pi$\;
    \BlankLine
    \tcp{3e. Update $\theta$ via L-BFGS-B or Adam}
    $\theta_{k+1} \leftarrow \theta_k + \alpha\, g$ \tcp*{or L-BFGS-B step}
    \BlankLine
    \tcp{3f. Dual stopping criteria (Gleave \& Toyer 2022)}
    \lIf{$\|g\| < \varepsilon_1$ \textbf{or} $\|D_{\text{data}} - D_\pi\|_\infty < \varepsilon_2$}{\textbf{break}}
}
\BlankLine

\tcp{Step 4: Inference}
$\ell(\hat\theta) = \sum_i \log \pi_{\hat\theta}(a_i | s_i)$\;
Bootstrap SEs by resampling individuals and re-estimating\;
\BlankLine

\KwOut{$\hat\theta$ (reward weights), $V^*$, $\pi^*$, log-likelihood, bootstrap SEs}
\end{algorithm}

\paragraph{Implementation notes.}
The backward pass (Step 3a) uses the hybrid solver from EconIRL, which starts with successive approximation and switches to Newton-Kantorovich near the fixed point for quadratic convergence. The forward pass (Step 3b) solves the occupancy measure via matrix inversion rather than iterative power method. The dual stopping criterion (Step 3f) follows \citet{gleave2022}, checking both gradient convergence and occupancy distance. L-BFGS-B is the default optimizer because it handles the implicit curvature well, but Adam is available for gradient-based optimization when the Hessian is ill-conditioned.

The deep variant replaces $\theta^\top \phi(s,a)$ with a neural network $r_\psi(s,a)$ following \citet{wulfmeier2016}. The gradient becomes $\nabla_\psi r$ weighted by the occupancy mismatch. This gains flexibility at the cost of interpretability and valid standard errors.


%% ============================================================
\section{Results}
%% ============================================================

% Auto-generated by mce_irl_run.py. Do not edit by hand.
% Known-truth DGP cells: canonical_low_action, mce_low_high_reward

\subsection{Known-Truth Validation Results}
The validation target is the tabular maximum causal entropy estimator with known transitions and known action-dependent reward features. Hard gates are enforced on feature matching, occupancy moments, normalized reward/value/Q recovery, policy distance, and Type A/B/C counterfactual regret.

\subsection{Validation Design Context}
\begin{itemize}
\item \textbf{low-dimensional action-dependent DGP.} A compact finite-state DDC benchmark with action-specific rewards, known transitions, known reward, value, policy, and Q-function truth, and Type A/B/C counterfactual oracles.
\item \textbf{low-state high-reward-feature DGP.} A maximum-causal-entropy stress cell with a compact state space and a richer action-dependent reward-feature basis.
\end{itemize}

\begin{table}[H]
\centering\small
\caption{MCE-IRL known-truth validation cells.}
\begin{tabular}{llrrrrr}
\toprule
Cell & Role & States & Actions & Reward features & Individuals & Periods \\
\midrule
low-dimensional action-dependent DGP & low-dimensional check & 21 & 3 & 4 & 2,000 & 80 \\
low-state high-reward-feature DGP & primary validation & 25 & 3 & 8 & 3,000 & 100 \\
\bottomrule
\end{tabular}
\end{table}

\begin{table}[H]
\centering\small
\caption{MCE-IRL primary fit summary.}
\begin{tabular}{lr}
\toprule
Quantity & Value \\
\midrule
Converged & true \\
Outer iterations & 25 \\
Log likelihood & -329244.2754 \\
Estimation time & 2.83 seconds \\
Optimizer & \texttt{root} \\
Feature residual & 0.000000 \\
Occupancy moment residual & 0.001060 \\
\bottomrule
\end{tabular}
\end{table}

\begin{table}[H]
\centering\small
\caption{MCE-IRL recovery metrics.}
\begin{tabular}{lrr}
\toprule
Metric & Sanity cell & Primary cell \\
\midrule
Feature residual & 0.000000 & 0.000000 \\
Occupancy moment residual & 0.000814 & 0.001060 \\
Reward normalized RMSE & 0.007252 & 0.082287 \\
Policy TV & 0.004529 & 0.006984 \\
Value normalized RMSE & 0.007398 & 0.082646 \\
Q normalized RMSE & 0.026314 & 0.081560 \\
Type A regret & 0.000126 & 0.000433 \\
Type B regret & 0.000516 & 0.000410 \\
Type C regret & 0.000139 & 0.000094 \\
\bottomrule
\end{tabular}
\end{table}

\begin{table}[H]
\centering\small
\caption{Hard-gate audit by cell.}
\begin{tabular}{lrr}
\toprule
Cell & Passed gates & Failed gates \\
\midrule
low-dimensional action-dependent DGP & 10 & 0 \\
low-state high-reward-feature DGP & 10 & 0 \\
\bottomrule
\end{tabular}
\end{table}

Raw parameter cosine is intentionally not a validation gate for MCE-IRL because feature matching, occupancy moments, reward/value/Q recovery, policy distance, and counterfactual regret are the decision-critical validation objects.


We run MCE-IRL, NFXP, and behavioral cloning on a \mceGridSize$\times$\mceGridSize\ gridworld with \mceNstates\ states, 5 actions, and three action-dependent features (step penalty, terminal reward, distance weight). Data consists of \mceNobs\ observations at $\beta = \mceBeta$.

MCE-IRL and NFXP achieve identical policy accuracy (\mcePolicyAcc\% and \mceNfxpPolicyAcc\%), confirming that they solve the same MLE. Both substantially outperform behavioral cloning (\mceBcPolicyAcc\%), which does not learn from MDP structure. The cosine similarities differ because rewards are identified only up to an additive constant and scale, so the raw parameter vectors may point in different directions while producing the same policy. Policy accuracy is the correct evaluation metric for IRL.

The feature matching gradient (Equation~\ref{eq:gradient}) is the key difference from NFXP. NFXP computes the likelihood gradient via implicit differentiation through the Bellman fixed point, which requires solving a linear system at each step. MCE-IRL computes the gradient as a difference between empirical and expected features, which requires the forward occupancy pass but avoids the implicit differentiation. Both approaches have the same computational cost per iteration (one backward pass and one forward solve), but MCE-IRL's gradient has a cleaner interpretation as feature matching.

\paragraph{Code.}
The companion script \texttt{mce\_irl\_results.py} in this folder generates Table~\ref{tab:mce_results}. The MCE-IRL estimator is implemented in \texttt{src/econirl/estimation/mce\_irl.py}. To regenerate:
\begin{verbatim}
python papers/econirl_package/primers/mce_irl/mce_irl_results.py
\end{verbatim}


\bibliographystyle{plainnat}
\bibliography{references}

\end{document}
